\chapter{Overview}

AI Code Reviewer reviews a code change with more than the diff. It parses the changed files, runs static rules, retrieves similar code from an index of the repository, and asks a language model for findings. Those findings are stored, shown in a dashboard, and, for a GitHub pull request, posted as one review with inline comments.

The same pipeline is available before a push. An MCP server exposes tools that Claude Code can call, including a review of a local \code{git diff}.

This is a university demo. The interesting part is the path from a pull-request event to a ranked set of comments, not multi-tenant scale. The stack runs locally with Docker Compose, or as a split deployment with the dashboard on Vercel and the API, worker, and database on Railway.

\section{What a review is}

A review is one run of the pipeline. It records who triggered it, which model answered, how many tokens it used, how long it took, a short summary, and a list of findings. Each finding has a file, a line range, a severity, a category, a message, and an optional suggestion.

Severity is one of critical, high, medium, low, or info. Category is one of bug, security, performance, style, or maintainability. A finding comes from the model or from Semgrep.

Two triggers matter in the demo:

\begin{itemize}
  \item A GitHub pull request is opened, updated, or reopened on a repository that has reviews turned on. Draft pull requests are ignored.
  \item A developer sends a raw diff through the MCP tool \code{review\_diff}.
\end{itemize}

A third trigger, manual re-run, repeats the review of a pull request from the dashboard. Eval runs use the same pipeline on a fixed set of diffs and do not post to GitHub.

\section{The five parts}

The repository is a pnpm workspace. Each part meets the others at the HTTP API and at the shared Zod schemas. Table~\ref{tab:modules} is the split of ownership.

\begin{table}[H]
  \centering
  \caption{Modules and what each one owns.}
  \label{tab:modules}
  \begin{tabular}{@{}lp{9.6cm}@{}}
    \toprule
    Module & Owns \\
    \midrule
    Git integration & GitHub App, webhooks, installation tokens, posting the review \\
    Indexer & Clone, Tree-sitter chunking, embeddings, pgvector search \\
    Review engine & Queue worker, prompts, Semgrep, validation and ranking, model calls \\
    Dashboard & Next.js app, Clerk sign-in, repository and review screens \\
    MCP and contracts & MCP tools and prompts, API keys, shared Zod schemas \\
    \bottomrule
  \end{tabular}
\end{table}

The four apps live in \code{apps/}: web, api, worker, and mcp. Shared libraries live in \code{packages/}: db (Drizzle schema and migrations), shared (Zod schemas, queue names, job payloads), and llm (one interface over several model providers). \code{evals/} holds the sample diffs used to score the reviewer.

\section{What the demo shows}

A walkthrough of the system is: sign in, create the GitHub App from the dashboard, enable a repository, index it, open a pull request that contains a known bug, watch the pipeline timeline, read the inline comments on GitHub, and run the same reviewer on a local diff from Claude Code. Chapter~\ref{chap:run} is the concrete runbook. The chapters in between describe how that walkthrough works.
